\section{Evaluasi}
\label{sec:evaluasi-bab5}

\subsection{Soal Pilihan Ganda}

\begin{enumerate}
    \item Setelah \lstasm|MOV AL, 20h| dijalankan pada register
          \reg{AX} yang semula berisi \hex{1234}, isi \reg{AX}
          menjadi \dots
          \begin{enumerate}
              \item \hex{1220} \item \hex{2034}
              \item \hex{0020} \item \hex{2020}
          \end{enumerate}

    \item Register yang berpasangan dengan \reg{CS} untuk menentukan
          alamat instruksi berikutnya adalah \dots
          \begin{enumerate}
              \item \reg{SP} \item \reg{IP}
              \item \reg{DS} \item \reg{DI}
          \end{enumerate}

    \item Instruksi \lstasm|ADD AL, BL| dengan \reg{AL} = \hex{7F}
          dan \reg{BL} = \hex{01} menghasilkan \dots
          \begin{enumerate}
              \item \flag{CF} = 1 dan \flag{OF} = 0
              \item \flag{CF} = 0 dan \flag{OF} = 1
              \item \flag{CF} = 1 dan \flag{OF} = 1
              \item \flag{CF} = 0 dan \flag{OF} = 0
          \end{enumerate}

    \item Register yang menunjuk puncak \eng{stack} adalah \dots
          \begin{enumerate}
              \item \reg{BP} \item \reg{SP}
              \item \reg{SI} \item \reg{DI}
          \end{enumerate}

    \item \eng{Flag} yang bernilai 1 ketika hasil operasi bernilai
          nol adalah \dots
          \begin{enumerate}
              \item \flag{CF} \item \flag{SF}
              \item \flag{ZF} \item \flag{OF}
          \end{enumerate}
\end{enumerate}

\subsection{Soal Uraian}

\begin{enumerate}
    \item Jelaskan mengapa \lstasm|MOV AL, 5| dan \lstasm|MOV AX, 5|
          menghasilkan keadaan register yang berbeda bila sebelumnya
          \reg{AH} tidak bernilai nol. Sertakan contoh angka.

    \item Sebutkan peran khas \reg{BX}, \reg{CX}, dan \reg{DX}.
          Jelaskan mengapa peran tersebut muncul walaupun keempat
          register umum sama-sama dapat menampung nilai.

    \item Diberikan keadaan awal \reg{AL} = \hex{80} dan
          \reg{BL} = \hex{80}. Tentukan hasil
          \lstasm|ADD AL, BL| beserta nilai \flag{CF}, \flag{OF},
          \flag{SF}, dan \flag{ZF}. Jelaskan mengapa \flag{CF} dan
          \flag{OF} bernilai demikian.

    \item Jelaskan peran \reg{CS}, \reg{DS}, \reg{SS}, dan \reg{ES}
          beserta pasangan alamatnya. Jelaskan pula akibat yang
          mungkin terjadi bila \reg{SS} dan \reg{SP} diubah secara
          sembarangan.
\end{enumerate}

\subsection{Kisi-kisi Penilaian}

\begin{table}[H]
    \centering
    \caption{Kisi-kisi evaluasi Bab~\ref{chap:register}.}
    \label{tab:kisi-bab5}
    \small
    \begin{tabular}{@{}p{1.5cm}p{5.0cm}p{2.6cm}c@{}}
        \toprule
        \textbf{Sub-CPMK} & \textbf{Indikator yang diukur} &
        \textbf{Bentuk soal} & \textbf{Bobot} \\
        \midrule
        1.3 & Mengidentifikasi register umum dan akses byte &
        Pilihan ganda 1; Uraian 1, 2 & 25\% \\
        \addlinespace
        1.3 & Menjelaskan register segmen dan pasangannya &
        Pilihan ganda 2; Uraian 4 & 25\% \\
        \addlinespace
        1.3 & Menjelaskan register \eng{pointer} dan \eng{index} &
        Pilihan ganda 4 & 15\% \\
        \addlinespace
        1.3 & Menganalisis perubahan \eng{flag} hasil operasi &
        Pilihan ganda 3, 5; Uraian 3 & 35\% \\
        \bottomrule
    \end{tabular}
\end{table}

\begin{checklistbox}
\vspace{1.1em}
\textbf{Checklist Pencapaian Kompetensi Bab~\ref{chap:register}}

Seorang mahasiswa dinyatakan kompeten pada Sub-CPMK 1.3 apabila:
\begin{itemize}
    \item[\checked] dapat menyebutkan register umum beserta byte
          tinggi dan byte rendahnya;
    \item[\checked] dapat menjelaskan peran khas AX, BX, CX, dan DX;
    \item[\checked] dapat menjelaskan peran CS, DS, SS, dan ES pada
          pengalamatan memori;
    \item[\checked] dapat menjelaskan tugas SP, BP, SI, DI, dan IP;
    \item[\checked] dapat menyebutkan sembilan \eng{flag} dan kapan
          masing-masing bernilai 1;
    \item[\checked] dapat memprediksi perubahan \eng{flag} pada
          operasi penjumlahan;
    \item[\checked] dapat menjelaskan pengaruh \eng{flag} terhadap
          keputusan program.
\end{itemize}
\end{checklistbox}
